% ===== sec/02_related_work.tex =====
\section{Related Work}
\label{sec:related}

The literature relevant to this paper falls into five strands.
The first asks whether agents are automata, and answers
affirmatively in a memory-architecture sense
(Koohestani et al.~\cite{koohestani2025automata}).  The second is
the producer--consumer framing of cooperating sequential
processes (Dijkstra~\cite{dijkstra1968cooperating}), which is
the operational form of the question we actually wish to answer:
how can the type of one agent bound another?  The third is the
classical theory of automata and the Chomsky hierarchy (Turing,
Chomsky, Rabin and Scott, Hopcroft and Ullman, Aho).  The fourth
is software verification and validation (Hoare, Dijkstra, Pnueli,
Clarke and Emerson, Lamport).  The fifth is the boundary
between decidable and undecidable verification fixed by Robinson
and the MRDP theorem.  We address each in turn.

\subsection{Are agents just automata?}

Koohestani, Li, Podkopaev, and Izadi~\cite{koohestani2025automata}
recently posed and answered the title question, establishing a
memory-architecture equivalence: simple-reflex agents are finite
automata; hierarchical task-decomposition agents are pushdown
automata; agents employing read/write memory for reflection are
Turing machines.  Their paper introduces a \emph{right-sizing
principle}: choose the minimal computational class sufficient
for the task.

Two observations are owed to that paper.  First, the equivalence
they establish is the architectural premise of the present work;
without it, the question of how Chomsky classes interact across
agents would have no canonical formulation.  Second, their paper
is preliminary in scope: it classifies agents but does not
prescribe a construction in which one agent's class
\emph{constrains} another's class.  We supply that construction
here.  Where they ask \emph{Are agents just automata?}\ and
answer \emph{yes}, we ask \emph{can the type of one agent bound
another?}\ and answer \emph{yes, with V\&V techniques inferred
into a Data-Driven verifier and the idealised theoretical
apparatus of automata built into a Goal-Driven producer}.

A related strand, Formal-LLM~\cite{li2024formalllm}, encodes
plan-construction constraints as a context-free grammar and
guides an LLM-based planner with a pushdown automaton.  Our
construction is dual:  we do not constrain plan generation
through a higher-class grammar, we \emph{verify} plan emission
through a lower-class one.  Aho--Corasick is the
``less-Turing-complete-than-the-LLM'' device that earns us the
verifier-side budget of Theorem~\ref{thm:closure}.

\subsection{The producer--consumer framing}

Dijkstra's cooperating sequential
processes~\cite{dijkstra1968cooperating} fix the producer--
consumer formulation:  one process produces a sequence of
artefacts; another consumes them; correct cooperation requires
that the consumer never reads an artefact the producer has not
yet produced, and the producer never overwrites an artefact the
consumer has not yet consumed.  The classical question is:
\emph{what synchronisation discipline is sufficient?}  In our
setting the question becomes: \emph{what Chomsky-class discipline
on the consumer is sufficient to render the producer's
recursively enumerable output decidable in finite time?}  The
answer, by Theorem~\ref{thm:closure}, is Type-3.  Lower would not
work (Type-3 is the minimal class that contains the keyword
matchers Aho--Corasick implements); higher would work but raise
the verifier-side budget unnecessarily and risk the halting
problem returning through the consumer's loop.

\subsection{Automata theory and the Chomsky hierarchy}

Turing's 1936 paper~\cite{turing1936computable} introduces the
abstract machine that bears his name and proves the existence of
a universal one.  Every formal claim in this paper inherits, at
some depth, from that construction.  Chomsky's 1956 and 1959
papers~\cite{chomsky1956three,chomsky1959certain} fix the
hierarchy itself: regular, context-free, context-sensitive, and
unrestricted (Type-3 through Type-0).  Rabin and
Scott~\cite{rabin1959finite} fix the formal correspondence
between regular languages and finite automata, with the
nondeterminism that makes the proofs tractable.  Hopcroft,
Motwani, and Ullman~\cite{hopcroft2006automata} consolidate the
modern treatment with strict inclusions and standard parsing
complexities.  Hartmanis and
Stearns~\cite{hartmanis1965computational} introduce the time and
space complexity classes in which our Type-1\,/\,Type-0 distinction
ultimately lives, and supply the formal vocabulary for our
``idealised Type-0 abstraction versus realised LBA'' discussion in
Section~\ref{sec:discussion}.

Aho's 1968 paper~\cite{aho1968indexed} introduces indexed
grammars as a strict extension of context-free grammars,
positioned strictly between Type-2 and Type-1.  In our framework
the back-edge structure of an iterative repair loop --- e.g.,
the \texttt{retry\_phase\_on\_low\_score} cycle in
Section~\ref{sec:methodology} --- inhabits exactly that zone.  We
use Aho's indexed grammars as the bridge between the pushdown
sub-agent (Type-2) and the linear-bounded sub-agent (Type-1).

The Aho--Corasick paper~\cite{aho1975corasick} compiles a finite
keyword set into a deterministic finite automaton that recognises
the union of the keywords in linear time and constant space.  In
our framework Aho--Corasick is \emph{the} canonical realisation
of the Data-Driven Type-3 consumer.

\subsection{Software verification and validation}

Hoare's axiomatic basis~\cite{hoare1969axiomatic} and Dijkstra's
guarded-command calculus~\cite{dijkstra1975guarded} supply the
machinery of Hoare triples $\{P\}\,C\,\{Q\}$ that we apply to
each unit-test step in the Aho--Corasick keyword set.  The
predicates $P$ and $Q$ are the ``before'' and ``after'' states
of a single tool invocation in the trace.

Pnueli's 1977 paper~\cite{pnueli1977temporal} introduces the
linear temporal logic in which we phrase the safety and
liveness properties of the Goal-Driven loop:
\[
G\bigl(\mathit{plan} \rightarrow F\,\mathit{verified}\bigr)
\]
expresses that every plan emitted by the producer is eventually
accepted or refuted by the consumer.  Clarke and
Emerson~\cite{clarke1981ctl} introduce the branching variant we
use to reason about the back-edges of the iterative repair loop,
and Clarke, Grumberg, and Peled~\cite{clarke1999modelchecking}
consolidate the model-checking framework that, in principle,
decides the Type-3 component of our verifier.

Lamport~\cite{lamport1978clocks,lamport1994tla} supplies the
ordering on events in the multi-agent execution and the temporal
logic of actions; the typesetting of this paper itself uses his
\LaTeX\ system~\cite{lamport1986latex}.  Liskov's abstract data
types~\cite{liskov1974abstraction} are the type discipline
underlying the \texttt{TypedDict} state objects that flow between
sub-agents.

\subsection{The undecidability boundary: Robinson and MRDP}

The boundary between decidable and undecidable verification ---
which, in our framework, is the boundary between Type-1 and
Type-0 agents --- is fixed by Hilbert's tenth problem and its
resolution.  Robinson's 1949 paper~\cite{robinson1949definability}
and the Davis--Putnam--Robinson follow-up~\cite{davis1961decision}
established that the recursively enumerable subsets of
$\mathbb{N}^k$ are exactly the Diophantine subsets, and that
the membership problem for an arbitrary recursively enumerable
set admits no general algorithm.  Our idealised Goal-Driven
producer is analysed against this recursive-enumeration backdrop,
but the empirical verifier does not decide full producer membership.
It decides the narrower and operationally useful question of whether
a candidate trace conforms to the regular verifier specification.

\subsection{Complexity, learning, and symbolic AI}

Cook~\cite{cook1971complexity} and Karp~\cite{karp1972reducibility}
fix the NP-completeness frame within which the model-search step
of the Goal-Driven producer (the Optuna tuning of LightGBM, the
cross-validated stacking ensemble) is, formally, an instance of
constrained combinatorial optimisation.
Yao~\cite{yao1977probabilistic} supplies the notion of expected
complexity used when reporting average trace lengths over the
input distribution.  Blum~\cite{blum1967machine} fixes the
machine-independent measure in which we report the validated
trace language's complexity.

The supervised-learning component of the Spaceship Titanic
solution is justified by the deep-learning manifesto of LeCun,
Bengio, and Hinton~\cite{hinton2015deep}; the causal discipline
that prevents the verifier from being fooled by spurious
correlations is justified by Pearl~\cite{pearl2009causality}.
McCarthy~\cite{mccarthy1960lisp} is the symbolic representation
of programs as data --- the property that lets a Goal-Driven
agent emit a \emph{program} (a Kaggle notebook) as the terminal
element of its trace.  Newell and Simon~\cite{newell1976physical}
state the physical-symbol-system hypothesis under which the
agent's tool calls are the elementary symbols.  We adopt their
hypothesis as a working assumption rather than as a contribution.

\subsection{Synthesis}

The strands above suffice to support every formal claim in the
sequel.  In particular, our novelty relative to
Koohestani et al.~\cite{koohestani2025automata} is the
producer--consumer construction: where they classify, we bound;
where they say which class an agent inhabits, we say which class
\emph{must} consume an agent of class $i$ in order for the system
as a whole to remain decidable.  The empirical instantiation in
Sections~\ref{sec:expsetup}--\ref{sec:results} is reported as
data, not as theory, and stands or falls on its own
reproducibility artefacts.
